% Generated draft: requires mathematical and source review.
\documentclass[11pt]{article}
\usepackage[margin=25mm]{geometry}
\usepackage{fontspec}
\setmainfont{Latin Modern Roman}
\usepackage{amsmath,amssymb,mathtools,unicode-math}
\setmathfont{Latin Modern Math}
\usepackage{enumitem,needspace,xurl,hyperref}
\setlength{\parindent}{0pt}
\setlength{\parskip}{5pt}
\setlength{\emergencystretch}{3em}
\newcommand{\N}{\mathbb{N}}
\newcommand{\Z}{\mathbb{Z}}
\newcommand{\Q}{\mathbb{Q}}
\newcommand{\R}{\mathbb{R}}
\newcommand{\C}{\mathbb{C}}
\newcommand{\F}{\mathbb{F}}
\newcommand{\A}{\mathbb{A}}
\newcommand{\PP}{\mathbb{P}}
\newcommand{\E}{\mathbb{E}}
\newcommand{\eps}{\varepsilon}
\newcommand{\dd}{\,\mathrm{d}}
\begin{document}
\section*{UnsolvedMath ID: 701353}
\textbf{Source problem code:} catalogue entry
\subsection{Definitions and mathematical statement}
Let \(G\) be a topological group with identity \(e\). Its topology is non-discrete if not every singleton is open. A subset \(D\subseteq G\) is dense if \(\overline{D}=G\). A partition of \(G\) into dense subsets is a family \((D_i)_{i\in I}\) of nonempty, pairwise disjoint dense subsets satisfying \(G=\bigcup_{i\in I}D_i\). Interpret countable as at most countable. A nontrivial subgroup \(H\leq G\) has exponent \(2\) if \(h^2=e\) for every \(h\in H\).

Assume that \(G\) is non-discrete and has no countable open subgroup of exponent \(2\). The supplied problem asks whether each of the following assertions holds for every such \(G\):
\begin{enumerate}
\item There exist dense subsets \(D_0,D_1\subseteq G\) such that
\[
G=D_0\sqcup D_1.
\]
\item There exist an infinite index set \(I\) and dense subsets \((D_i)_{i\in I}\) such that
\[
G=\bigsqcup_{i\in I}D_i.
\]
\end{enumerate}
The parenthetical wording ``two (infinitely many)'' is retained as these two questions; no prescribed infinite cardinality is supplied. The record reports that the answer is affirmative when \(G\) is Abelian or countable, under the same hypotheses.

The source record does not specify a separation axiom for topological groups. Thus the intended class of topological groups, in particular whether Hausdorffness is implicit, remains an unresolved qualification of this formulation.
\subsection{Short English statement}
Must every non-discrete topological group with no countable open subgroup of exponent \(2\) admit a partition into two dense subsets, and also a partition into infinitely many dense subsets? The source reports affirmative answers for Abelian or countable groups. Its separation-axiom convention requires clarification.
\subsection{Sources}
\begin{itemize}
\item Local UnsolvedMath registry, permanent ID \(701353\), exact source problem code \texttt{catalogue entry}, title \emph{Kourovka Notebook Problem 17.51}. The supplied registry labels the problem open; this status has not been independently verified.
\item \emph{The Kourovka Notebook (Unsolved Problems in Group Theory)}, original problem \(17.51\), 17th issue (2010); the supplied background also identifies the 21st issue edition (2026), \texttt{arXiv:1401.0300v48}. These bibliographic details are supplied attribution and have not been independently checked.
\item Proposed by E. G. Zelenyuk and I. V. Protasov, with proposed year 2010 in the supplied record. The collection description names E. I. Khukhro and V. D. Mazurov as editors.
\end{itemize}
\end{document}
